% =============================================================================
% Project : THE_PHILOSOPHY_OF_VIABILITY
% Book    : Book I — Ontology of Viability
% Book ID : VS.B01
%
% File    : 4_ASSETS/diagrams/VS.B01.P01.C02.S11.SS10.P004.DG01.dimensions_of_an_invariance_claim.tex
%
% Asset   : DG01 — Dimensions of an Invariance Claim
% Type    : Diagram
% Location: Section 11, Subsection 10, Paragraph 004
%
% Purpose :
%   To represent the distinction and coordination of claim
%   specification, form of preservation, and domain of examination.
%
% Requires (preamble):
%   \usepackage{tikz}
%   \usetikzlibrary{arrows.meta, positioning, fit}
% =============================================================================

\begin{figure}[htbp]
    \centering
    \hyphenpenalty=10000 \exhyphenpenalty=10000\relax  % no hyphenation inside diagram labels

    \begin{tikzpicture}[
        font=\small,
        node distance=0.4cm,
        % --- base box ---
        box/.style={
            draw, rounded corners=3pt, align=center,
            inner sep=6pt, line width=0.7pt
        },
        % --- node roles (names avoid clashes with built-in TikZ keys) ---
        dimbox/.style={box, text width=3.3cm, minimum height=4.1cm,
                       fill=blue!8, draw=blue!55!black},
        assessbox/.style={box, text width=5.4cm, minimum height=1.1cm,
                          font=\small\bfseries,
                          fill=blue!16, draw=blue!55!black},
        resultbox/.style={box, text width=5.4cm, minimum height=1.1cm,
                          dashed, fill=green!8, draw=green!40!black},
        % --- frames and labels ---
        vsgroup/.style={draw=black!45, dashed, rounded corners=5pt,
                        inner sep=8pt, line width=0.5pt},
        grouplabel/.style={font=\footnotesize\itshape, text=black!75,
                           align=center},
        % --- arrows ---
        directed/.style={-{Stealth[length=2.5mm]}, line width=0.9pt}
    ]

    % ---------------------------------------------------------------
    % Three analytical dimensions (equal rank, no order among them)
    % ---------------------------------------------------------------
    \node[dimbox] (specification) {
        \textbf{Claim Specification}\\[5pt]
        \footnotesize
        Subject, determination, transformation,
        respect, criterion, conditions, scope
    };
    \node[dimbox, right=of specification] (form) {
        \textbf{Form of Preservation}\\[5pt]
        \footnotesize
        Determinational, structural, relational, compositional,
        quantitative, qualitative, configurational and combined forms
    };
    \node[dimbox, right=of form] (domaindim) {
        \textbf{Domain of Examination}\\[5pt]
        \footnotesize
        Four proposed fundamental domains
        and three intermediate hubs
    };

    \node[vsgroup, fit=(specification)(form)(domaindim)] (dims) {};
    \node[grouplabel, above=0.1cm of dims] {
        Distinct analytical dimensions, jointly coordinated
    };

    % ---------------------------------------------------------------
    % Assessment and result
    % ---------------------------------------------------------------
    \node[assessbox, below=0.9cm of dims] (assessment) {
        Invariance Assessment
    };
    \node[resultbox, below=0.9cm of assessment] (result) {
        Preservation claim within a specified scope
    };

    \draw[directed] (dims)       -- (assessment);
    \draw[directed] (assessment) -- (result);

    % ---------------------------------------------------------------
    % Note
    % ---------------------------------------------------------------
    \node[below=0.7cm of result, text width=10cm, align=center,
          font=\footnotesize\itshape] {
        The assessment does not independently establish
        numerical identity, persistence, continuity,
        or the total ontological constitution of the existent.
    };

    \end{tikzpicture}

    \caption{The three analytical dimensions of an invariance
    claim. They coordinate the assessment of preservation
    without becoming interchangeable or establishing conclusions
    beyond the specified scope.}
    \label{fig:dimensions-of-an-invariance-claim}
\end{figure}

% =============================================================================
% END OF DIAGRAM DG01
% =============================================================================
